La géométrie est le lieu où l'on apprend à démontrer. C'est là que le collège introduit la
différence entre constater sur une figure et établir par un raisonnement, et cette
différence est le vrai contenu du thème.

Les chapitres suivent le mouvement qui va de la figure au calcul. On commence par les objets
et leurs propriétés — triangles, quadrilatères, cercle, Thalès, Pythagore — puis par les
transformations qui les déplacent sans les déformer, puis par les grandeurs qu'on leur
attache. Le lycée change d'outil sans changer de sujet : le vecteur remplace la figure, le
produit scalaire remplace le rapporteur, et les mêmes théorèmes se redémontrent en une ligne
de calcul — Al-Kashi n'est que le développement du carré d'une différence.

Deux avertissements valent pour tout le thème. Les démonstrations ici ne s'appuient jamais
sur une figure : les points sont des vecteurs, et rien ne peut se déduire d'un dessin bien
orienté. Et les aires et volumes forment le point dur du collège, parce que la formule de
l'aire du disque, donnée à l'école, est en mathématiques un théorème d'intégration ; le
chapitre concerné le dit à sa place.
